\documentclass[a4paper,11pt]{article}

\usepackage{times}
\usepackage[T1]{fontenc}
%\usepackage[xetex]{graphicx} % when you compile with xelatex
\usepackage[pdftex]{graphicx} % when you compile with pdflatex
\usepackage{epstopdf}
\usepackage{float}
\usepackage{enumerate}
\usepackage{boxedminipage}
\usepackage{url}
%\usepackage[dutch]{babel} % when you write in dutch

\usepackage{dcolumn}
\newcolumntype{d}[1]{D{.}{.}{#1}}

%%%%%%%%%%%% layout %%%%%%%%%%%%%%%%%%
\usepackage[margin=3cm]{geometry}
% you might enhance readability by increasing the distances between lines:
%\renewcommand{\baselinestretch}{1.1}

%%%%%%%%%%% mathematics %%%%%%%%%%%%%%%%%%
\usepackage{amsmath,amsfonts,amssymb}

%%%%%%%%%%%%% theorems %%%%%%%%%%%%%%%%
% adapt when you write Dutch

\usepackage{amsthm}

\theoremstyle{plain}
\newtheorem{theorem}{Theorem}
\newtheorem{corollary}{Corollary}
\newtheorem{proposition}{Proposition}
\newtheorem{lemma}{Lemma}
\newtheorem{assumption}{Assumption}

\theoremstyle{definition}
\newtheorem{definition}{Definition}
\newtheorem{example}{Example}
\newtheorem{remark}{Remark}
\newtheorem{notation}{Notation}

%%%%%%%%%%%%algorithms %%%%%%%%%%%%%%%%%%%
\usepackage{algorithm}
\usepackage{algorithmic}

%%%%%%%%% bibliography style %%%%%%%%%%%%%
\usepackage{natbib} % in case of classic bibtex
%\usepackage[natbib,style=authoryear,maxcitenames=2,
%           maxbibnames=99,hyperref,backend=biber]{biblatex}
%\addbibresource{bib/EORthesis.bib}

%%%%%%%%%%%%%%%% start document %%%%%%%%%%%
\begin{document}

%%%%%%%%%%%%%%% front page %%%%%%%%%%%%%%%%

\thispagestyle{empty}

\includegraphics[height=2cm]{figures/LogoSBE.png}

\vspace*{3cm}

\noindent
\rule{\textwidth}{0.8pt}
\begin{center}
{\huge\bf
\noindent
Deriving fractional moments using the Moment Generating Function
}
\end{center}

\vspace*{-8pt}
\noindent
\rule{\textwidth}{0.8pt}

\vspace*{2cm}

\begin{center}
{\LARGE\bf
Jelle Reisinger
}

{\Large
\vspace*{0.5cm}
(2780350)


\vspace*{2cm}

30 June 2025
}
\end{center}

\vspace*{2cm}

{\Large
\noindent
Bachelor Thesis Econometrics and Data Science
}

\vspace*{1cm}

{\Large
\noindent
Thesis commission:\\[0.3cm]
Prof. dr. xxxx (supervisor)\\[0.3cm]
Dr. yyyy (co-reader)
}


\newpage

%%%%%%%%%%%%%%%%%%%%%%% main contents %%%%%%%%%%%%%%%%%
\setcounter{page}{1}

\section*{Abstract}
The abstract should summarize the contents of the thesis.
It should be clear, descriptive, self-explanatory and not longer
than a third of a page. Please avoid using mathematical
formulas as much as possible.
Keywords might be given.

\bigskip\noindent
\textbf{Keywords:} Forecasting, Optimization.

\section{Introduction}\label{s:intro}
We prefer and recommend that you use LaTeX for preparing your
bachelor thesis.
This paper contains a set of instructions to get the required formats and styles.
Please follow the guidelines herein when preparing your thesis.
The easiest way to write your thesis using LaTeX that complies with the
requirements is to edit the source file, \texttt{template-and-guide.tex,}
for this document. In case you choose for an alternative document preparation system
(for instance Microsoft Word), be sure to obtain the same formats and styles.
And make sure that you can export your document to pdf format.

\section{General Guidelines}\label{s:general}
\subsection{Article Style}
The thesis is prepared in LaTeX using an article style, e.g., standard
LaTeX article style

\verb+\documentclass[a4paper,11pt]{article}+.

\noindent
The source file (the \texttt{.tex} version) is typeset using \texttt{pdflatex},
or \texttt{xelatex} which allows you to use graphics file formats 
\texttt{pdf, jpg, png, eps}
by including the \texttt{graphicx} package.

\subsection{Language}
We prefer that you prepare your thesis in English (UK or US), but we allow Dutch.
Please carefully check the spelling of words before you submit your thesis.
There are spell checkers for LaTeX as well.
Some examples of software which supports spell checking are
TexnicCenter, TexMaker, TeXShop, and WinEdt.

\subsection{Font Specification and Spacing}
The thesis should be set in the Times New Roman font using a 11-point font size.
These settings are obtained by the option 11pt (in \texttt{documentclass})
and by the package \texttt{times}.
The thesis should be single spaced which you may increase slightly by a factor 1.1:

\verb+\renewcommand{\baselinestretch}{1.1}+.

\subsection{Margins}
The thesis is printed on A4 paper.
The text area is 15.0 cm wide and the 23.7 cm high, centered, which means that
all margins are 3 cm. This is simply obtained by

\verb+\usepackage[margin=3cm]{geometry}+.

\subsection{Footnotes}
Do not use footnotes; instead incorporate such material into the text
directly or parenthetically.

\subsection{Page Numbers}
The pages are numbered except the front page.

\subsection{Formatting the Front Page}
The front page just contains
\begin{itemize}
\setlength\itemsep{-1mm}
\item
Title of the thesis.
\item
Author (your name and student number).
\item
Date.
\item
References or logo(s) of
\begin{itemize}
\item
VU University.
\item
School of Business and Economics.
\item
Department of Econometrics and Operations Research.
\end{itemize}
\item
Master Thesis.
\item
Thesis Committee and its members.
\end{itemize}


\section{Mathematics}
\subsection{Mathematical Expressions in Text and in Displays}
Display only the most important equations, and number only the
displayed equations that are explicitly referenced in the text.
To conserve space, simple mathematical expressions such as
$\bar Y = n^{-1} \sum_{i=1}^n Y_i$ may be incorporated into the text.
Mathematical expressions that are more complicated or that must be
referenced later should be displayed, as in
\[
s^2 = \frac 1 {n-1} \sum_{i=1}^n (Y_i - \bar Y)^2.
\]

If a display is referenced in the text, then enclose the equation number
in parentheses and place it flush with the right-hand margin of the
column. This is automatically obtained by the \texttt{equation} environment
accompanied by the \verb+\label+ command.
For example, the quadratic equation has the general form

\begin{equation} \label{eq:quadratic}
ax^2 + bx + c = 0, \mbox{ where } a \ne 0.
\end{equation}

In the text, each reference to an equation number should also be enclosed in
parentheses by using \verb+\eqref{<labelgiven>}+.
For example, the solution to \eqref{eq:quadratic} is given
in \eqref{eq:quadraticsol} in Appendix \ref{app:quadratic}.

If the equation is at the end of a sentence, then you should end the equation
with a period. If the sentence in question continues beyond the equation,
then you should end the equation with the appropriate
punctuation---that is, a comma, semicolon, or no punctuation mark.

\subsection{Symbols, Commands, Environments, Etc}
See

\url{https://en.wikibooks.org/wiki/LaTeX/Mathematics}

\noindent
and

\url{https://en.wikibooks.org/wiki/LaTeX/Advanced_Mathematics}.


\subsection{Definitions, Theorems, Etc.}
Definitions, theorems, propositions, etc. should be formatted
using the \texttt{amsthm} package.
Number these items separately and sequentially. Examples are given below

\begin{definition}
In colloquial New Zealand English, the term \textit{dopey mongrel} is used
to refer to someone who has exhibited less than stellar intelligence.
\end{definition}

\begin{theorem}\label{t:delete}
If a proceedings editor from New Zealand accidentally deletes his draft of
the author kit shortly after completing it, he would be considered to be a dopey mongrel.
\end{theorem}
\begin{proof}
The proof follows by the principle of contradiction.
Suppose the editor is not a dopey mongrel, then he is smart enough to save
the author kit.
\end{proof}

\begin{corollary}
One of the proceedings editors is a dopey mongrel.
\end{corollary}
\begin{proof}
This follows immediately from Theorem \ref{t:delete}.
\end{proof}

\section{Figures and Tables}
\label{sec:graphics}
Figures and tables should be centered within the text and should not
extend beyond the right and left margins of the paper.
Figures and tables can make use of color.
However, try to select colors that can be differentiated when printing
in black and white in consideration of vast majority of people
using such printers.
Figures and tables are numbered sequentially, but separately,
using arabic numerals.

\subsection{Tables}
Each table should appear in the document after the paragraph in which the
table is first referenced. However, if the table is getting split across
pages, it is okay to include it after a few paragraphs from its first reference.
One-line captions are centered, while multiline captions are left justified.
The captions appear \textit{above} the table. See Tables \ref{tab:first} and
\ref{tab:second} for examples.

\begin{table}[H]
\caption{Table captions appear above the table, and if they are longer
than one line they are left justified. Captions are written using normal
sentences with full punctuation. It is fine to have multiple-sentence
captions that help to explain the table.}
\label{tab:first}
\medskip
\centering
\begin{tabular}{rll}
\hline
Creature & IQ & Diet\\ \hline
dog & 70 & anything\\
cat & 75 & almost nothing\\
human & 60 & ice cream \\
dolphin & 120 & fish fillet\\
\hline
\end{tabular}
\end{table}

\begin{table}[H]
\centering
\caption{Counting in Maori.}
\label{tab:second}
\medskip
\begin{tabular}{r|l}
English & Maori \\ \hline
one & tahi \\
two & rua \\
three & toru \\
four & wha \\
\hline
\end{tabular}
\end{table}

\begin{table}[H]
  \centering
  \caption{Alternative table aligning columns}
  \label{tab:alt}
  \medskip
  \begin{tabular}{r|d3d5}
       & \multicolumn{1}{c}{2012} & \multicolumn{1}{c}{2013} \\
    \hline
    Mean & 1.23 & 1.23456 \\
    Variance & 13454.4 & 3435.456 \\
    \hline
  \end{tabular}

  \medskip
  \parbox{100mm}{\small Note: Alternatively, place a longer explanation
  as a note at the bottom of the table, allowing the caption to be more
  concise. This table uses explicit alignment of the numbers in the
  columns, using either 3 or 5 decimals.}
\end{table}


\subsection{Figures}
Each figure should appear in the document after the paragraph in which
the figure is first referenced. One-line captions are centered,
while multiline captions are left justified. Figure captions appear
below the figure. 
To include figures you use package \texttt{graphicx}, and in the document the
command

\verb+\includegraphics{<graphs/graphicfilename>}+.

The graphicfilename usually does not have to include an extension;
PdfLaTeX would search for extensions it recognises. In general, it is
advisable to keep your graphs in a separate directory, to avoid clogging
up your LaTeX directory.

If the graph is scaled correctly, one should be able to use a

\verb+\includegraphics[width=\textwidth]{<graphs/graphicfilename>}+

to ensure that the graph fills the width of the paper.




\noindent



\subsection{References to Tables and Figures}
References to tables and figures identified by number are capitalized.
For example, ``We see in Table \ref{tab:second} that...'' and
``We see in the previous table that...'' are both correct.
Be sure to use the \verb+\label+
command within the figure or table environment and refer to the
associated figure or table using \verb+Table \ref{<labelgiven>}+.
Please do not use hard coded figure/table numbers.
This is error prone.

\subsection{Graphics Formats}
As graphics files in your document you use \texttt{.jpg}, \texttt{.png}, 
\texttt{.pdf}, or \texttt{.eps} files.
But there are tools to convert these formats into one another.
The main difference between the formats is how they store the images and how
well suited they are for specific graphics. In general we can
choose between bitmap and vector graphics.
Bitmap graphics are well suited for photographies (jpg is very common here)
or for screenshots (png is a lossless encoding (in contrast to jpg),
and is thus better suited for all those cases where you have sharp edges in your graphics).
Vector graphics are the encoding to be chosen for all kinds of drawings
(diagrams, charts, ...). In contrast to bitmap formats, they can be scaled
to any size without any loss of sharpness. This makes it possible to
read such graphics even if two pages are printed on one sheet of paper,
or if the documents are read electronically.

So what to choose for your Latex document? As a rule of thumb you should
always prefer pdf and eps. In general these two encodings can contain both,
bitmap and vector graphics. But there is no need (and no use) to convert
your bitmaps to any of these.

You include figures via the \verb+\includegraphics+ command.
You must use the \texttt{pdflatex} or \texttt{xelatex} command to generate
your pdf file, as was done with this file.

\section{Algorithms}
Typeset algorithms by using the \texttt{algorithmic} environment of the
\texttt{algorithm} package.
The command \verb+\begin{algorithmic}+ can be given the optional argument of a
positive integer, which if given will cause line numbering to occur at multiples
of that integer. E.g. \verb+\begin{algorithmic}[5]+ will enter the algorithmic
environment and number every fifth line. Below is an example of typesetting a
basic algorithm (remember to add the

\verb+\usepackage{algorithm,algorithmic}+

\noindent
statement to your document preamble). The pseudocode is centered by using the
\texttt{minipage} environment.

{\small
\begin{verbatim}
\begin{center}
\begin{minipage}{10cm}
\begin{algorithm}[H]
\caption{Polar Method for Normal Random Numbers}
\begin{algorithmic}[1]
\REPEAT
   \STATE $U_1\sim {\sf U}(0,1)$ \COMMENT{Generate uniform on $(0,1)$}
   \STATE $U_2\sim {\sf U}(0,1)$
   \STATE $V_1 \gets 2U_1 -1$ \COMMENT{Uniform on $(-1,1)$}
   \STATE $V_2 \gets 2U_2 -1$
   \STATE $W \gets V_1^2 + V_2^2$
\UNTIL {$W < 1$}
\STATE \textbf{return} $V_1\,\sqrt{(-2\,\ln(W)/W}$
\end{algorithmic}
\end{algorithm}
\end{minipage}
\end{center}
\end{verbatim}
}

\noindent
This produces

\begin{center}
\begin{minipage}{10cm}
\begin{algorithm}[H]
\caption{Polar Method for Normal Random Numbers}
\begin{algorithmic}[1]
\REPEAT
   \STATE $U_1\sim {\sf U}(0,1)$ \COMMENT{Generate uniform on $(0,1)$}
   \STATE $U_2\sim {\sf U}(0,1)$
   \STATE $V_1 \gets 2U_1 -1$ \COMMENT{Uniform on $(-1,1)$}
   \STATE $V_2 \gets 2U_2 -1$
   \STATE $W \gets V_1^2 + V_2^2$
\UNTIL {$W < 1$}
\STATE \textbf{return} $V_1\,\sqrt{(-2\,\ln(W)/W}$
\end{algorithmic}
\end{algorithm}
\end{minipage}
\end{center}

\section{Bibliography Management}
Managing your bibliography requires two specifications:
the citation style in your document,
and the style of the reference list at the end of the thesis.

\subsection{Citing a Reference}
To cite a reference in the text, use the author-date method. Thus,
\citet{ross06} could also be cited parenthetically \citep{ross06}.
For a work with three or more authors, use an abbreviated form.
For example, a work by Evans, Keith and Kroese would be cited in
one of the following ways:
\citet{evkekro07} or \citep{evkekro07}.

Parenthetical citations are enclosed in parentheses $(~)$, not square brackets $[~]$.
The items in a series of such citations are usually separated by commas.
If an item in the series of parenthetical citations contains punctuation 
because (for example) it refers to a work with three or more coauthors, 
then all items should be separated by semicolons.

The following is a list of correct forms of citations:
\begin{itemize}
\setlength\itemsep{-1mm}
\item Brown and Edwards (1993),
\item (Brown and Edwards 1993),
\item (Brown and Edwards, 1993),
\item Brown and Edwards (1993), Smith (1997),
\item (Brown and Edwards 1993; Smith 1997; Brown et al. 1997).
\item (Brown and Edwards, 1993; Smith, 1997; Brown et al., 1997).
\end{itemize}

\subsection{List of References}
Place the list of references after the appendices.
The section heading is \textbf{References}, and it is not numbered.
List only references that are cited in the text.
Arrange the references in alphabetical order (chronologically for a
particular author or group of authors); do not number the references.
Give complete references without abbreviations.
To identify multiple references by the same authors and year, append a
lower case letter to the year of publication; for example, 1984a and 1984b.

Use hanging indentation to distinguish individual entries.
Do not insert additional space between references.

The bibliographic style for a journal article is: \\
$<$Surname of first author$>$, $<$First author initials$>$,
$<$Initials and surnames of other authors$>$ ($<$year$>$).
$<$Article title$>$. $<${\em Journal Name in
Headline Italics}$>$ $<$Volume number$>$, $<$page numbers$>$.

The article title may come between quotes but this is not required.
The format for other types of reference can be inferred from the
examples in the references, which include:
\begin{itemize}
\setlength\itemsep{-1mm}
\item a technical report \citep{chi89},
\item a proceedings article \citep{evkekro07},
\item a journal article \citep{alon},
\item a book by 2 authors \citep{asmus07},
\item a chapter in a book \citep{asrub1},
\item an unpublished thesis or dissertation \citep{garvels00},
\item a document available on the web \citep{sudoku}.
\end{itemize}

\subsection{BibTeX}
These formats are obtained by creating a BibTeX database for all your references.
A BibTeX database is stored as a \texttt{.bib} file.
It is a plain text file, and so can be viewed and edited easily.
One benefit of using BibTeX is that bibliography formatting and
referencing can be greatly simplified: the correct citation and
reference list style is automatically created.
How the BibTeX items are entered in the \texttt{bib} file will be
explained in section \ref{ss:bibtex}.

Once you have created the
\texttt{bib} file with your references it needs to be included
in your document.

\subsection{Biblatex Package}
Advise is to use the \texttt{biblatex} package. It is a reimplementation of the
classic BibTeX providing modern bibliographic facilities such as
advanced name disambiguation, smart crossref data inheritance,
configurable sorting schemes, dynamic datasource modification, etc.
The prerequisite is that you need to use Biber to process your \texttt{bib}
files. Older TeX distributions do not have Biber included, you need to check this
or you find out when you cannot run your \texttt{bib} files using the
\texttt{biblatex} package.
In that case, see next section for the classic BibTeX implementation.

\begin{itemize}
\item
In the preamble you put

\begin{verbatim}
\usepackage[natbib,
            style=authoryear,
            maxcitenames=2,
            maxbibnames=99,
            hyperref,
            backend=biber]{biblatex}
\addbibresource{bib/<bibfilename>.bib}
\end{verbatim}

\item
At the end of your document, where the reference list is supposed to come you
put

\verb+\printbibliography+

\item
In your document, wherever you need
to cite a reference you put \verb+\citet{<bibentrykey>}+ (for the author name(s)
followed by the year in parentheses), or
\verb+\citep{<bibentrykey>}+
to get the entire citation in parentheses.
\end{itemize}

\subsection{Classic BibTeX}
In case Biber does not work in your TeX distribution.
Use the \texttt{natbib} bibliography style with \texttt{chicago} citation style.

\begin{itemize}
\item
In the preamble you put

\verb+\usepackage{natbib}+.

\item
At the end of your document, where the reference list is supposed to come you
put

\begin{verbatim}
\bibliographystyle{chicago}
\bibliography{bib/<bibfilename>.bib}
\end{verbatim}

\item
Citations are \verb+\citet{<bibentrykey>}+ (for the author name(s)
followed by the year in parentheses), or
\verb+\citep{<bibentrykey>}+
to get the entire citation in parentheses.
\end{itemize}

\subsection{BibTeX Entries}\label{ss:bibtex}
See the \texttt{EORthesis.bib} file available in Canvas,
or use the templates that you can find, e.g.,

\url{https://en.wikibooks.org/wiki/LaTeX/Bibliography_Management}

\noindent
and

\url{https://en.wikipedia.org/wiki/BibTeX}.

\subsection{Compilation}
You need to obtain the final pdf version of document with correct citations and reference
list.

\begin{enumerate}[(a).]
\item
In case you use Biblatex, you run consecutively
\texttt{pdflatex, biber, pdflatex}. In stead of
\texttt{pdflatex}, you might also run \texttt{xelatex}.
\item
In case you use classic BibTeX, you run consecutively
\texttt{pdflatex, bibtex, pdflatex, pdflatex}. In stead of
\texttt{pdflatex}, you might also run \texttt{xelatex}.
\end{enumerate}

\subsection{Example}
In \citet{asmus07} we see that they mention \citep{blitzstein10},
and after some more comments of \citet{asrub1}, we know that
\citet{evkekro07} have published the same results
as we can find in \citep{alon,garvels00,ross06},
however in a more general context.

\section{Conclusion}\label{s:con}
The conclusion summarizes the results of your research.
Also it gives possible directions of further research.

\section*{Acknowledgements}
Place the acknowledgments section, if needed, after the main text, 
but before any appendices and the references. The section heading is not numbered.

\newpage
\appendix

\section{Appendix}\label{s:app1}
Place any appendices after the acknowledgments, starting on a new page.
The appendices are numbered
\textbf{A}, \textbf{B}, \textbf{C}, and so forth.
The appendices contain material that you would like to share with the
reader but that would hinder the flow of reading. For instance long
proofs of theorems, code of algorithms, data, etc.

\section{Appendix}\label{s:app2}
You might give more appendices. For instance Appendix \ref{s:app1} for
proofs, Appendix \ref{s:app2} for data.

\section{Appendix}\label{app:quadratic}
The solution to \eqref{eq:quadratic} has the form

\begin{equation} \label{eq:quadraticsol}
x = \frac{-b \pm \sqrt{b^2-4ac}}{2a} \quad\text{if}\;\; a \ne 0.
\end{equation}

%\printbibliography
% or in case of classic bibtex:
\bibliographystyle{chicago}
\bibliography{bib/bibfile}

\end{document}
